\chapter[User Interface Implementation \textsuperscript{Novotny}]{User Interface Implementation}
\label{chap:UI}
The implementation of the user interface is based on QML. \\
This chapter explains how the user interface was constructed. QML and the components that have been used for building the user interface are explained in chapter 3.
\paragraph{Organisation of the Chapter} The first part of this chapter gives an overview of the initialization process required before the first menu page is loaded. The second part explains the architecture of the user interface. The next section explains how the different methods of menu navigation have been implemented. The third section deals with the components used throughout the menu pages. At the end of the chapter a number of sections deal with side processes (such as storing and retrieving user data) that are triggered by the user interface.
\section{Initialization and Data Supply}
An {\tt ApplicationViewer} is responsible for displaying the user interface.
In order to allow communication between the \cpp and the QML side of the application two context properties are passed to the {\tt ApplicationViewer} object:

\begin{itemize}
	\item{{\tt licenseManager}}
	\item{{\tt game}}
\end{itemize}

The {\tt game} property points to the {\tt MainGame} object instantiated in the main method (refer to \ref{chap:Design}) for an explanation of the role of this object). This property is the most important communication interface between the user interface and the application logic.
Its main responsibilites include:
\begin {itemize}
	\item {Supplying data such as {\tt Scores}, {\tt Tracks} and {\tt Devices}.}
	\item {Controlling the flow of the application by triggering the start of a song and maintaining state information (refer to the chapter \textit{Design} for a more detailed explanation).}
\end {itemize}

A {\tt SoundNavigationHandler} object responsible for passing sound-based navigation instructions which has a pointer to the {\tt ApplicationViewer} object is created right after the user interface is initialized. The responsibilites of this object are discussed in detail at a later point.

\section {Menu Structure}
When the application starts a QML component called {\tt DisplayLoader} is loaded. 
Its responsibilites include 
\begin {itemize}
	\item {Dynamically loading the menu pages. This is accomplished using the {\tt Loader Element}. When {\tt DisplayLoader} is first created the Main Menu is loaded. Thereafter, the individual menus are responsible for resetting the source of the item to be displayed in the DisplayLoader.
	\begin{lstlisting}
		Loader {
			id: loader
			anchors.fill: parent
		}
	\end{lstlisting}
	The menus may change the source of the loader as follows.
	\begin{lstlisting}
		CustomButton {
			id: startButton
			buttonText: "Start"
				
			onClicked: {
			     	loader.source = "GameSettings.qml"
			 }
		}
	\end{lstlisting}
	}
	\item {
	Supplying global parameters concerning the environment and the appearance of the application. These include:
	\begin {itemize}
		\item{appearance parameters such as font, font-size and color}
		\item{environment parameters such as name and location of the database, API links, ...}
		\item{layout parameters such as insets and margins}
	\end {itemize}
	}
\end {itemize}
\section{Navigation}
As explained in chapter \ref{chap:UserInteractionConcept} the application supports three user interaction concepts.
\begin{itemize}
	\item {Mouse-Based Navigation}
	\item {Sound-Based Navigation}
	\item {Keyboard-Based Navigation}
\end{itemize}
\subsection{Sound-Based Navigation}
\label{sec:sound based navigation}
Commonly used navigation commands can be triggered by playing chords on a connected instruments. The following conditions need to be satisfied in order for a chord to be recognised:
\begin {itemize}
	\item {Chords consist of notes rather than pitches.}
	\item {Chords have to consist of three distinct notes.}
	\item {All notes of the chord need to be audible at the same time.}
	\item {Different inversions of the same chord are not distinguished. A \textit{D/F \#} chord is handled the same way a \textit{D} chord is handled.}
\end{itemize}
Playing a chord  while satifsying these conditions will be referred to as issuing a \textit{Sound-Based navigation instruction}.
\paragraph{Workflow}
This paragraph gives an overview of the process that is triggered in consequence of a single Sound-Based navigation instruction. The steps will be described in detail below.
\begin {enumerate}
	\item {The {\tt SoundNavigationHandler} recognizes a chord.}
	\item {It invokes the QML function {\tt handleCommand} defined in the {\tt DisplayLoader} element.}
	\item {The function emits a signal of the currently loaded menu page.}
	\item {This signal is received by the {\tt Chord} element of the current menu page. If a {\tt Chord} element finds out that its notes have been played it emits the signal {\tt played}.}
	\item {The signal is handled by the menu page.}
\end {enumerate}
\paragraph{Recognizing Chords}
Recognising chords is not necessary during game situations as single notes contained in the chord are evaluated separately. Chord recognition has been implemented for the sole purpose of allowing Sound-Based navigation. Chord recognition is done by the class {\tt SoundNavigationHandler}. {\tt SoundNavigationHandler} inherits {\tt QObject}. An object of this class is instantiated as soon as the {\tt ApplicationViewer} that displays the graphical user interface has been created. The {\tt ApplicationViewer} is the parent of the {\tt SoundNavigationHandler} object which means it will be destroyed automatically once the graphical user interface is destroyed. The class offers two slots:
\begin {itemize}
	\item{{\tt noteOn}}
	\item{{\tt noteOff}}
\end {itemize}
These slots are connected to the {\tt noteOn} and {\tt noteOff} signals of the {\tt Input} class while the menus are being used similarly to how these signals are connected to the {\tt Evaluator} class during game situations. \\
The object stores all notes that are currently being played (which means a {\tt NoteOn} but not corresponding {\tt NoteOff} event has been received) in a collection. As soon as there are three notes being played at the same time it considers the combination a chord. Note that only chords containing exactly three notes can be handled. 
The {\tt SoundNavigationHandler} object has a pointer to the {\tt ApplicationViewer} object. It can therefore invoke QML functions. This is used for communicating with the graphical user interface. 
\paragraph{Emitting the signal {\tt chordPlayed}}
The function {\tt handleCommand} that is defined in the {\tt DisplayLoader} element is invoked by the {\tt SoundNavigationHandler}. It receives a list containing the notes in the chord. It determines which menu page is currently loaded and then triggers the signal {\tt chordPlayed} of the active menu page. The signal is defined in the {\tt NavigatablePanel} element which is the root element of all menu pages.
\paragraph{Chord Element} The chord element is an invisible QML item. It stores an array of notes which make up the chord it represents. When a chord element is created it connects the parent's signal {\tt chordPlayed} to its function {\tt chordPlayed} which acts as a slot. This function decides whether or not the chord that was played is the chord represented by the element. It receives three midi pitch values. At this point midi pitch values are converted to musical notes represented by the letters A to G, using the \#-sign to indicate a pitch rise by a semi-tone. This is done using a modulo operation since musical notes repeat every twelve semitones. As a result the {\tt Chord} element discards the octave of the notes. This means that any inversion of a chord will be recognized.
\paragraph{Handling the {\tt Played} signal of a chord}
Six {\tt Chord} elements which are responsible for the most common user inputs are created in the {\tt NavigatablePanel} element. These include:
\begin {itemize}
	\item {Four chords corresponding to the arrow keys on the keyboard. If one of these chords is played the handling is done by the {\tt NavigatablePanel} element in a generic way that applies to all menu pages. This will be described in detail at a later point.}
	\item {A chord responsible for selecting the currently highlighted item of a {\tt SelectableListView} is defined by {\tt SelectableListView}}.
	\item {Additionally, two chords representing the following Commands are defined by {\tt NavigatablePanel}: 
		\begin {itemize}
			\item {instant confirmation of the settings made and proceeding to the next step}
			\item {discarding all settings and returning to the last step}
		\end {itemize} 
	A corresponding signal that can be handled by the active menu page is emitted when one of these chords is played.}
\end {itemize}
\subsection{Keyboard-Based Navigation}
When a key is pressed an event is delivered to the QML item that is currently focused. If the focused item does not accept the event it is passed up to its parent item until the event is handled or the root element has been reached. \\
Handling key events is done at different spots inside the user interface:
\begin {itemize}
	\item {
		Arrow key- and tab key events are handled by the root element of the active menu page. The root element of the menu pages is of type {\tt NavigatablePanel}.
	}
	\item {
		Return key events are handled by custom components such as {\tt CustomButton} and {\tt CustomCheckBox}. Using the return key when one of these components is focused has the same effect as clicking on the component.
	}
	\item {
		Space key events cause selection of the currently highlighted item of a {\tt SelectableListView}.
	}
\end {itemize}
The following code example shows how key events are handled by the {\tt CustomButton} element. 
\begin{lstlisting}
     // Keyboard Handling
     Keys.onPressed: {
         if (event.key === Qt.Key_Return && button.focus) {
             event.accepted = true;
             button.clicked();
         }
     }
\end{lstlisting}
Note that the event has to be accepted. Omitting this step will cause the event to be forwarded to the parent element of the button. Some buttons are used to change the source of a loader element which might cause the parent of the button to be destroyed which might lead to a segmentation fault.

\subsection{Generic Key Handling}
{\tt NavigatablePanel} handles a number of key events relevant to all inheriting menu pages:
\begin {itemize}
	\item {Arrow keys for moving inside a focused {\tt SelectableListView} or moving to the next component of the menu page. This is described in greater detail in the next section.}
	\item {The Enter key causes confirmation of all settings of the current dialog and proceeding to the next one. {\tt NavigatablePanel} emits the signal {\tt confirmed} which can be handled by the inheriting menu page.}
	\item {The Space key causes going back to the last dialog. The signal {\tt back} is emitted by {\tt NavigatablePanel} and may be handled by the inheriting menu page.}
\end {itemize}

\subsection{Generic Chord Handling}
{\tt NavigatablePanel} contains chord elements triggering certain actions: 
\begin {itemize}
	\item {one {\tt Chord} instance corresponding to each arrow key}
	\item {a {\tt Chord} instance corresponding to the enter key}
	\item {a {\tt Chord} instance corresponding to the space key}
\end {itemize}
The {\tt onPlayed} signals of these {\tt Chord} components are handled similarly to how the key events are handled.
\subsection{Focus Handling and Pre-Selection of List Items}
Focus Handling has been implemented in a generic way in the {\tt NavigatablePanel} element which is the root element of every menu page. \\
Every {\tt NavigatablePanel} element maintains an array called {\tt focusList} which contains all child elements that are part of the linear navigation flow. These elements are manually added to the element in the desired order when a menu page is created.
{\tt NavigatablePanel} provides the following functions.
\begin {itemize}
	\item {
		The function {\tt shiftFocus} triggers a focus shift by the passed amount. This means that the focus is given to another element of the {\tt focusList}. The parameter {\tt amount} can be positive (causing forward shifts) or negative (causing backward shifts). When the last element of the focus currently has the focus a positive shift will take the focus back to the beginning of the list. The opposite case is handled analogously.
	}
	\item {
		The function {\tt moveInActiveList} triggers a pre-selection shift in the {\tt SelectableListView} that currently has focus. It can be used similarly to {\tt shiftFocus}. When the last element of the list is selected a positive shift will not have any effect as the user might want to scroll to the end of the list quickly without being taken back to the beginning. The opposite case is handled analogously.
	}
\end {itemize}
The {\tt NavigatablePanel} has a property providing access to the currently focused component.

\subsection{Visual Feedback}
\subsubsection{Focused Component}
All custom components that can be focused ({\tt SelectableLisView}, {\tt CustomButton}, {\tt CustomCheckbox} ...) extend {\tt HighlightablePanel}.{}
\paragraph{\tt HighlightablePanel} consists of two {\tt Image} components:
\begin{itemize}
	\item {{\tt BackgroundImage}}
	\item {{\tt SourceImage}}
\end{itemize}
The {\tt source} property of these components can be accessed via the {\tt imageSource} and {\tt highlightSource} property aliases of {\tt HighlightablePanel}. The inheriting component is responsible for setting the visibility of the images.
\section {Components}
\subsection{Standard Components, Reasons for Implementing Custom Components}
QML provides a number of basic components which serve as building blocks for custom components rather than functional parts of the user interface.
This type of components include:
\begin{itemize}
	\item{{\tt Item}}
	\item{{\tt Rectangle}}
	\item{{\tt MouseArea}}
\end{itemize}
and many others. These components are treated in Chapter 3.

The package {\tt QtQuick Controls} offers a range of controls commonly used in applications. The appearance of these components is specific to the platform and therefore provides a native look and feel. However, there is no way to customize the appearance of the components while keeping the logic. All user controls were reimplemented to suit the look and feel of the application using only basic QML components as building blocks.

\subsection{NavigatablePanel}
The role of the {\tt NavigatablePanel} component is providing functionalities required by all menu pages. 
All menu pages inherit {\tt NavigatablePanel}. 

The responsibilites of {\tt NavigatablePanel} include:
\begin {itemize}
	\item {storing and retrieving data to/from the database,}
	\item {generic key event handling for key-based navigation,}
	\item {generic chord handling for sound-based navigation,}
	\item {generic focus handling.}
\end {itemize}

The functionalites of {\tt NavigatablePanel} are described in more detail in the corresponding sections.

\subsection{CustomButton}
The {\tt CustomButton} implements the function of a button while meeting the optical requirements of the application's user interface.  Primarily, this means to display different background images depending on the state of the button (default, hovered, highlighted), and to enlargen the button text when the button is highlighted.  The main building block for this custom component is the {\tt MouseArea} component provided by QML. Apart from reacting to mouse clicks the component reacts to keyboard events. 


\includegraphics[width=0.25\textwidth]{custom_button_a}

\includegraphics[width=0.25\textwidth]{custom_button_b}

{\tt CustomButton} offers the signal {\tt clicked} which can be emitted from within the component, which is done when the {\tt MouseArea} component detects a click or a key event is handled, or from without the component which is required for sound-based navigation (refer to the section on sound-based navigation for a in-depth explanation).

The component offers the following properties:
\begin {itemize}
	\item {{\tt buttonText}}
	\item {{\tt fontColor} sets the font color to be used for displaying the button text. If the property is not set explicitly the default font color defined in a global will be used.}
	\item {{\tt fontSize} sets the font size to be used for displaying the button text. If the property is not set explicitly the default font size defined in a global will be used.}
	\item {{\tt imageSource} sets the source of the default background image}
	\item {{\tt mouseHighlightSource} sets the source of the background image to be displayed while the component is hovered or focused.}
\end {itemize}

\subsection{CustomCheckBox}
The {\tt CustomCheckBox} implements the functionality of check box. 
{\tt CustomCheckBox} inherits {\tt HighlightableItem}.
The space key can be used for checking and unchecking the check box.


\includegraphics[width=0.5\textwidth]{custom_checkbox}

It provides the following properties:
\begin {itemize}
	\item {{\tt text} sets the text to be displayed next to the check box}
	\item {{\tt checked} indicates whether or not the check box is checked}
	\item {{\tt checkedSource}/{\tt uncheckedSource} set the source of the graphic to be displayed when check box is checked or unchecked}
\end {itemize}

\subsection{CustomSlider}
{\tt CustomSlider} implements the functionality of selecting a numeric value within a specified range while preventing the user from selecting illegal values. {\tt CustomSlider} components have a default value.

\includegraphics[width=0.5\textwidth]{custom_slider}

{\tt CustomSlider} inherits {\tt HighlightablePanel}.
Resetting the slider to the default value can be achieved by double clicking on the slider.
Apart from using mouse-based navigation to change the selection, the following key commands have been implemented:

\begin {itemize}
	\item {Horizontal arrow keys for changing the selection}
	\item {Space key for resetting the selection to the default value}
\end {itemize}

It offers the following properties:
\begin {itemize}
	\item {{\tt min} and {\tt max} set the lowest and highest value respectively}
	\item {{\tt defaultValue} contains the default value of the slider. Unless reassigned, the value of the property is set to the average of the {\tt min} and {\tt  max} properties.}
	\item {{\tt value} contains the currently selected numeric value}
\end {itemize}

\subsection{SelectableListView}
{\tt SelectableListVew} is a {\tt View} component that displays data supplied by a model. It allows pre-selection and selection of items using mouse, keyboard and Sound-Based navigation. The implementation is based on the QML {\tt ListView} component.
{\tt SelectableListView} defines the appearance and behaviour of the container element. By default {\tt SelectableListView} uses a {\tt SelectableListViewDelegate} and assumes the model provides scalar data. 

The component offers the follwing properties: 
\begin {itemize}
	\item {
		{\tt itemDelegate} allows passing customised delegates to the component. This allows displaying complex data in a more sophisticated way. In order for the component to work properly, {\tt itemDelegate} needs to be a {\tt SelectableListViewDelegate}.
	}
	\item {
		{\tt model} defines the data to be displayed by the list view. This data is made available to the delegate and can therefore be arbitrarly complex.
	}
	\item {
		{\tt selectedIndex}
	}
	\item {
		{\tt preSelectedIndex}
	}
	\item {
		{\tt selectedItem} provides access to the currently selected data item.
	}
	\item {
		{\tt lastItem} points to the last item that has been selected. This is required for triggering a exitinguishing effect on the delegate.
	}
	\item {
		{\tt itemWidth} sets the width of the delegates and determines how many items are displayed at once. This property becomes irrelevant if the {\tt itemDelegate} property is set manually.
	}
	\item {
		{\tt text} sets the header of the component
	}
\end {itemize} 

\paragraph{Highlighting Items}
A highlight effect is applied to the currently selected item of a {\tt SelectableListView}. Highlighting of selected items of a {\tt ListView} has been implemented in QML. Highlighting and transitions have, however, been reimplemented to better suit the look and feel of the application. An illuminating and distinguishing effect when changing the selection has been implemented. When an item is selected, {\tt SelectableListView} needs to keep track of both, the index of the item that has been selected and the index of the item that has been deselected. It will then instruct the corresponding delegates to activate or deactive highlighting.

\subsection{SelectableListViewDelegate}
The generic list view delegate called {\tt SelectableListViewDelegate} has been developed to meet the visual and functional requirements of most of the list views throughout the application. It is highly flexible in terms of the data it displays and can be used in the following ways.
\paragraph{Properties and Ways of Using the Delegate}
\begin {itemize}
	\item {
		\textbf{Displaying scalar data}
		If no optional properties of the delegate are set it assumes a scalar data model. This used to display single-valued data like the different instrumental tracks available in a song.
	}
	\item {
		\textbf{Displaying customised text}
		The property {\tt itemText} can be set manually. This can be used to define how compound data is displayed.
	}
	\item {
		\textbf{Manually defining the data item}
		By manually setting the property {\tt dataItemSource} the way the delegate displays data can be overriden. This way, look and feel as well as functionality of the {\tt SelectableListViewDelegate} can be maintained while providing complete flexibility in terms of contents and layout.
	}
\end {itemize}
The following excerpt of the implementation of {\tt SelectableListViewDelegate} is sufficient for explaining how the three ways of using the component have been implemented.
\begin{lstlisting}
Item {
    property string itemText: (typeof modelData === typeof "" ? modelData : "")
    property alias dataItemSource : dataItemLoader.sourceComponent

    // Access data of current item from outside delegate
    property variant dataOfItem : model

    Item {
        id: standardDataItem
        anchors.fill: parent
        z: 1

        Text {
            anchors.top: parent.top
            anchors.left: parent.left
            anchors.margins: marginInItem

            z: 1
            text: itemText
            color: standardFontColor
            font.pixelSize: standardHeaderPixelSize
            font.bold: true
            elide: Text.ElideRight
            width: parent.width  - 2 * marginInItem
        }

        visible: dataItemLoader.source !== undefined
    }

    Loader {
        id: dataItemLoader
        anchors.fill: parent
        z: 2
    }
}
\end{lstlisting}
The property {\tt itemText} determines what text is displayed by the delegate. Its default value is {\tt modelData} which contains the value of an associated scalar record. The property can be redefined explicitly in order to allow displaying compound data in a specific way. In order to avoid errors caused by assigning compound model data to the property, the type of the value is checked. \\
The property {\tt dataItemSource} is relevant for the third way of using the delegate. If the property is set an {\tt Item} defined by the user will be loaded dynamically and displayed instead of the text. \\
Note how the property {\tt visible} of the text element is bound to the {\tt source}-property of the {\tt Loader}.

The property {\tt dataOfItem} allow accessing the data associated with the delegate from outside the delegate. 
\paragraph{Highlighting}
The delegate provides the functions {\tt highlightOn()} and {\tt highlightOff()}. The corresponding {\tt SelectableListView} component is responsible for calling the functions on the right delegates.

\subsection{Search}
In order to facilitate finding the desired score file, a search function has been implemented.
\subsubsection{Search component}
The {\tt Search} component is intended to be used as a child of a {\tt SelectableListView} component. It will call its parent's {\tt applyFilter} function. The {\tt Search} component offers a property {\tt text} which provides access to the filter text and a signal {\tt filterTextChange} to notify the user about changes to the filter text.
\subsubsection{Score Search}
A {\tt Search} component is placed inside the {\tt SelectableListView} responsible for displaying the available score files.
\paragraph{Filtering}
The function {\tt applyFilter} is implemented as a child of the {\tt SelectableListView} instance. The function {\tt applyFilter} updates the data model of the {\tt SelectableListView}. 
\begin{lstlisting}
	function applyFilter(text) {
	    model = mainGame.getAvailableScoresFiltered(text)
	}
\end{lstlisting}

The \cpp function {\tt getAvailableScoresFiltered} returns a {\tt QList} of {\tt Score-} objects.
A score object is considered a match to the filter text if 
\begin {itemize}
	\item {the filter text is contained in the song title}
	\item {the filter text is contained in the name of the artist}
\end {itemize}.
\begin{lstlisting}
	QList<QObject *> MainGame::getAvailableScoresFiltered(QString filterText)
	{
	    QList<QObject*> scores;
	
	    for (int i = 0; i < availableScores.length(); i++) {
	        if (availableScores.at(i)->getTitle().contains(filterText) || availableScores.at(i)->getArtist().contains(filterText)) {
	            scores.append((QObject*) availableScores.at(i));
	        }
	    }
	
	    return scores;
	}
\end{lstlisting}
\paragraph{Visibility}
The {\tt Search} component is only visible if a filter text has been set. If text is typed while the {\tt SelectableListView} the {\tt Search} component is associated with, the {\tt Search} component needs to become visible. If the filter text is deleted the component needs to become invisible.  
\paragraph{Focus and Key Forwarding}
If text is entered while the {\tt SelectableListView} component is focused, the focus is given to the {\tt Search} component. By the time the first letter is typed the {\tt SelectableListView} component still has the focus. This means that the {\tt Search} component misses out on the first letter that is typed. The first letter needs to be \textit{forwarded} to the {\tt Search} component. The {\tt SelectableListView} component handles key events and decides whether forwarding the input is necessary. If the input is forwarded focus is given to the {\tt Search} component.
\begin{lstlisting}
		// Keyboard Handling
		Keys.onPressed: {
		    if (event.key === Qt.Key_Return || event.key === Qt.Key_up) {
		        scoreListView.forceActiveFocus()
		        scoreSearch.text = scoreSearch.text.trim()
		
		    }
		
		    else if (event.text !== " ") {
		        scoreSearch.text += event.text
		        scoreSearch.forceActiveFocus()
		    }
		}
\end{lstlisting}
\section{Storing and Retrieving User Preferences}
\label{sec:Preferences}
In order to reduce the amount of necessary adjustments prior to starting the game certain settings are saved when changed and retrieved every time the application is started. 

A local SQLite Database is used to store the user preferences.

\paragraph{Qt Quick Local Storage}
The module Qt Quick Local Storage assists in storing and retrieveing data from SQLite databases. 
It handles automatic creation and connection to the local database.

\paragraph{Storing Data}
Data is stored as soon as all settings made inside the dialog haven been confirmed (e.g. when starting a game in the Game Start Menu). Inside the dialogs that require the data, a connection to the database is established. The dialogs are responsible for handling the transaction (and therefore splitting the instructions into separate transactions). The following code example shows how the functions defined in the {\tt NavigatablePanel} element are used to load user preferences inside the Game Settings dialog. The data is written to the database using the {\tt storePreferences()} function.

\begin{lstlisting}
function storePreferences() {
    var db = LocalStorage.openDatabaseSync(db_name);
    db.transaction(
        function(tx) {
            writePreferenceToDB(tx, 'Device', deviceListView.currentItem.itemText)
            writePreferenceToDB(tx, 'Display Option', displayOptionsListView.currentItem.itemText)
            writePreferenceToDB(tx, 'Training Mode', trainingModeCheckBox.checked)
        }
    )
}
\end{lstlisting}

The function {\tt LocalStorage.openDatabaseSync()} automatically creates the database if it does not exist. Due to small amount of records stored in the table truncating and refilling the table every time the data changes is acceptable.
\paragraph{Loading Data}
Data is loaded when ever a dialog requiring the data is loaded.
In order to avoid duplicating code the code that handles database queries has been implemented inside the {\tt NavigatablePanel} custom element which is inherited by all dialogs of the user interface.

The following code example shows how user preferences are loaded inside the Game Settings Dialog. The function {\tt loadPreferences()} fetches the stored user preference data from the database and uses the data to preset the controls of the dialog. The function {\tt indexOf()} finds the corresponding item in the {\tt ListView}s data model. If no corresponding item can be found the first item in the {\tt ListView} will be selected.
\begin{lstlisting}
    function loadPreferences() {

        var db = LocalStorage.openDatabaseSync(db_name)
        var displayOption
        var device
        var trainingMode

        db.transaction(
            function(tx) {
                // Query Database
                displayOption = askDB(tx, 'Display Option')
                device = askDB(tx, 'Device')
                trainingMode = askDB(tx, 'Training Mode')
            }
        )

        displayOptionsListView.currentIndex = indexOf(displayOptionsListView, displayOption)
        deviceListView.currentIndex = indexOf(deviceListView, device)
        trainingModeCheckBox.checked = trainingMode === 1 ? true : false
    }
\end{lstlisting}
Data required at multiple points throughout the application is loaded every time the main menu loaded. Creating the database and intializing the data is also done by the main menu if needed. 
\paragraph{Database Layout}
The table {\tt Preferences} consists of two columns (key, value) and stores key value pairs defining the user preferences. \\
The database does not imply any constraints on the data it keeps. The application is responsible for ensuring the data it stores is valid. In particular, this means ensuring that the keys stored in the table are unique. The records of the table are altered when the user preferences change. The only situation where new records are created is when the application is started for the first time or when the database file is removed.
\section {Displaying Album Art and Graphics}
\subsection {Displaying Album Art}
In order to facilitate finding the desired score and enhance the appearance of the application corresponding album art is displayed where possible. 

\includegraphics[width=\textwidth]{album_covers}

\paragraph{Workflow}
This is an overview of how album art for the individual score files is loaded.
\begin {enumerate}
	\item {A delegate representing a score file is created. At this point meta data of the score have been read from the file header and are requested by the delegate.}
	\item {The delegate calls the function {\tt requestImage} passing its corresponding song title and artists as well as a reference to delegate itself to the function.}	
	\item {The function {\tt requestImage}} calls the function {\tt request} (as explained below) passing the title of the song, the reference to the delegate as well as a callback function defined at this point.
	\item {The function {\tt request} asynchronously executes a HTTP reqeust to the iTunes API. The callback function that was defined in the last step will be called once the HTTP request is completed. The reference to the delegate and the content of the HTTP response will be passed to the callback function. }
	\item {The callback function parses the response and extracts the URL of the image to be displayed by the delegate.}
	\item {Finally the callback function calls a function of the delegate which will display the image located at the specified URL.}
\end {enumerate}
\paragraph{iTunes API}
The iTunes API supplies meta data on music releases and artists. Comminucation with the API is based on HTTP.
JSON is used for representing the delivered data.
The terms of use allow using promotional contents including song previews and album art for promotional purposes only.
\paragraph{HTTP communication in JavaScript}
The XMLHttpRequest API is used for sending HTTP requests in JavaScript.
The event listener {\tt onreadystatechange} will be invoked on four different occasions. The state of the request is stored in the {\tt readyState} property of the request object.
Only the last state change is of interest to the application. This occurs when the content of the response has finished loading. 
The method {\tt open} initializes the {\tt XMLHttpRequest}-object.
At this point the HTTP request method is specified. The third parameter causes the request to be executed asynchronously.
No data is sent to the server along with this request.
\begin{lstlisting}
    function request(url, del, callback) {
            var xhr = new XMLHttpRequest();
            xhr.onreadystatechange = (function(myxhr) {

                return function() {
                    // finished
                    if (xhr.readyState === 4) {
                        callback(myxhr, del);
                    }
                }

            })(xhr);
            xhr.open('GET', url, true);
            xhr.send('');
    }
\end{lstlisting}
\paragraph{Parsing JSON in JavaScript}
JSON is an open data format primarily intended for exchanging data between applications. JSON is a proper subset of JavaScript. For this reason parsing JSON in JavaScript turns out to be particularly neat.
\begin{lstlisting}
	// translate response into object
	var receivedData = eval('new Object(' + o.responseText + ')');
	
	// access elements inside json object with dot notation
	if (receivedData.results !== undefined) {
	    if (receivedData.results[0].artworkUrl100 !== undefined) {
	        var uri = receivedData.results[0].artworkUrl100
	        delegate.setImage(uri)
	    }
	}
\end{lstlisting}
The function {\tt eval()} evaluates or executes the supplied JavaScript code. In the listing above the function returns a new object containing the JSON-notated data coming rom the HTTP response. \\ Based on the structure of the JSON content accessing the data is straightforward.